% Options for packages loaded elsewhere
\PassOptionsToPackage{unicode}{hyperref}
\PassOptionsToPackage{hyphens}{url}
\PassOptionsToPackage{dvipsnames,svgnames,x11names}{xcolor}
\documentclass[
  french,
  10pt,
  a4paper,
]{article}
\usepackage{xcolor}
\usepackage[margin=2.2cm]{geometry}
\usepackage{amsmath,amssymb}
\setcounter{secnumdepth}{-\maxdimen} % remove section numbering
\usepackage{iftex}
\ifPDFTeX
  \usepackage[T1]{fontenc}
  \usepackage[utf8]{inputenc}
  \usepackage{textcomp} % provide euro and other symbols
\else % if luatex or xetex
  \usepackage{unicode-math} % this also loads fontspec
  \defaultfontfeatures{Scale=MatchLowercase}
  \defaultfontfeatures[\rmfamily]{Ligatures=TeX,Scale=1}
\fi
\usepackage{lmodern}
\ifPDFTeX\else
  % xetex/luatex font selection
\fi
% Use upquote if available, for straight quotes in verbatim environments
\IfFileExists{upquote.sty}{\usepackage{upquote}}{}
\IfFileExists{microtype.sty}{% use microtype if available
  \usepackage[]{microtype}
  \UseMicrotypeSet[protrusion]{basicmath} % disable protrusion for tt fonts
}{}
\makeatletter
\@ifundefined{KOMAClassName}{% if non-KOMA class
  \IfFileExists{parskip.sty}{%
    \usepackage{parskip}
  }{% else
    \setlength{\parindent}{0pt}
    \setlength{\parskip}{6pt plus 2pt minus 1pt}}
}{% if KOMA class
  \KOMAoptions{parskip=half}}
\makeatother
\usepackage{longtable,booktabs,array}
\usepackage{calc} % for calculating minipage widths
% Correct order of tables after \paragraph or \subparagraph
\usepackage{etoolbox}
\makeatletter
\patchcmd\longtable{\par}{\if@noskipsec\mbox{}\fi\par}{}{}
\makeatother
% Allow footnotes in longtable head/foot
\IfFileExists{footnotehyper.sty}{\usepackage{footnotehyper}}{\usepackage{footnote}}
\makesavenoteenv{longtable}
\ifLuaTeX
\usepackage[bidi=basic,shorthands=off,]{babel}
\else
\usepackage[bidi=default,shorthands=off,]{babel}
\fi
\ifLuaTeX
  \usepackage{selnolig} % disable illegal ligatures
\fi
\setlength{\emergencystretch}{3em} % prevent overfull lines
\providecommand{\tightlist}{%
  \setlength{\itemsep}{0pt}\setlength{\parskip}{0pt}}
% Fichier LaTeX généré par docs/build/build_pdf.sh à partir de 02_methode_conception.md.
% Compilation (depuis le dossier de ce fichier), deux fois pour la table des matières :
%   pdflatex 03_Guide2_Methode_de_conception.tex
% Pour insérer une capture : déposez le fichier PNG au nom indiqué dans le cadre « CAPTURE À INSÉRER »
% (dossier ../images/...), puis recompilez. Vous pouvez aussi remplacer la commande \CredixCapture{...}{...}
% par votre propre \includegraphics.
\newcommand{\CredixShortTitle}{Guide 2 : méthode}
% Préambule commun aux PDF du dossier de remise CREDIX (compilé avec pdfLaTeX via pandoc)
\usepackage[scaled=.92]{helvet}
\renewcommand{\familydefault}{\sfdefault}
\usepackage{pifont}
\usepackage{amssymb}
% Caractères Unicode utilisés dans les documents (pdfLaTeX ne les connaît pas d'office)
\DeclareUnicodeCharacter{03C1}{\ensuremath{\rho}}
\DeclareUnicodeCharacter{2265}{\ensuremath{\geq}}
\DeclareUnicodeCharacter{2264}{\ensuremath{\leq}}
\DeclareUnicodeCharacter{2248}{\ensuremath{\approx}}
\DeclareUnicodeCharacter{2192}{\ensuremath{\rightarrow}}
\DeclareUnicodeCharacter{2190}{\ensuremath{\leftarrow}}
\DeclareUnicodeCharacter{2212}{\ensuremath{-}}
\DeclareUnicodeCharacter{2713}{\ding{51}}
\DeclareUnicodeCharacter{2714}{\ding{52}}
\DeclareUnicodeCharacter{2717}{\ding{55}}
\DeclareUnicodeCharacter{25B2}{\ensuremath{\blacktriangle}}
\DeclareUnicodeCharacter{25CB}{\ensuremath{\circ}}
\DeclareUnicodeCharacter{2500}{-}
\DeclareUnicodeCharacter{2502}{|}
\DeclareUnicodeCharacter{250C}{+}
\DeclareUnicodeCharacter{2510}{+}
\DeclareUnicodeCharacter{2514}{+}
\DeclareUnicodeCharacter{2518}{+}
\DeclareUnicodeCharacter{251C}{+}
\DeclareUnicodeCharacter{2524}{+}
\DeclareUnicodeCharacter{252C}{+}
\DeclareUnicodeCharacter{2534}{+}
\DeclareUnicodeCharacter{253C}{+}
\DeclareUnicodeCharacter{0192}{\textit{f}}
\DeclareUnicodeCharacter{2011}{-}
\DeclareUnicodeCharacter{202F}{\,}
\DeclareUnicodeCharacter{2009}{\,}
\DeclareUnicodeCharacter{00D7}{\ensuremath{\times}}
\DeclareUnicodeCharacter{25A0}{\ensuremath{\blacksquare}}
\DeclareUnicodeCharacter{2022}{\ensuremath{\bullet}}
\DeclareUnicodeCharacter{2026}{\ldots}
\DeclareUnicodeCharacter{2460}{\ding{172}}
\DeclareUnicodeCharacter{2461}{\ding{173}}
\DeclareUnicodeCharacter{2462}{\ding{174}}
\DeclareUnicodeCharacter{2463}{\ding{175}}
\DeclareUnicodeCharacter{2464}{\ding{176}}
\DeclareUnicodeCharacter{2465}{\ding{177}}
\DeclareUnicodeCharacter{2466}{\ding{178}}
\DeclareUnicodeCharacter{2467}{\ding{179}}
\usepackage{fvextra}
\DefineVerbatimEnvironment{verbatim}{Verbatim}{breaklines=true,breakanywhere=true,fontsize=\footnotesize,frame=single,framerule=0.3pt,rulecolor=\color[gray]{0.7},framesep=2mm,xleftmargin=1mm,xrightmargin=1mm}
\fvset{breaklines=true,breakanywhere=true}
\usepackage{fancyhdr}
\usepackage{tcolorbox}
\usepackage{colortbl}
\usepackage{enumitem}
\definecolor{credixblue}{HTML}{1F3A5F}
\definecolor{credixgray}{HTML}{5F6B7A}
\definecolor{boxbg}{HTML}{F3F5F8}
\setlength{\emergencystretch}{4em}
\renewcommand{\arraystretch}{1.3}
\setlength{\parindent}{0pt}
\setlength{\parskip}{0.55em}
\setlist{itemsep=0.15em,topsep=0.3em}
\pagestyle{fancy}
\fancyhf{}
\fancyhead[L]{\small\color{credixgray}\textsc{\CredixShortTitle}}
\fancyhead[R]{\small\color{credixgray}CREDIX --- Dossier de remise}
\fancyfoot[C]{\small\color{credixgray}\thepage}
\renewcommand{\headrulewidth}{0.3pt}
\renewcommand{\footrulewidth}{0pt}
\fancypagestyle{plain}{\fancyhf{}\fancyfoot[C]{\small\color{credixgray}\thepage}\renewcommand{\headrulewidth}{0pt}}
\usepackage{titlesec}
\titleformat{\section}{\Large\bfseries\color{credixblue}}{}{0em}{}[\vspace{-0.2em}\color{credixblue}\titlerule]
\titleformat{\subsection}{\large\bfseries\color{credixblue}}{}{0em}{}
\titleformat{\subsubsection}{\normalsize\bfseries\color{credixblue}}{}{0em}{}
% Cadre de repère pour une capture manquante
\newcommand{\CaptureManquante}[2]{%
  \begin{tcolorbox}[colback=boxbg,colframe=credixgray,boxrule=0.5pt,arc=1mm,left=3mm,right=3mm,top=3mm,bottom=3mm,halign=center]
  \textbf{CAPTURE À INSÉRER}\\[0.4em]#1\\[0.4em]{\footnotesize\ttfamily #2}
  \end{tcolorbox}}

% Images : insérées si le fichier existe au moment de la compilation, sinon cadre « CAPTURE À INSÉRER ».
% Pour ajouter une capture : déposer le PNG au nom indiqué dans le cadre, puis recompiler.
\usepackage{graphicx}
\newcommand{\CredixCapture}[2]{%
  \IfFileExists{#2}{%
    \begin{center}\includegraphics[width=\linewidth,height=0.75\textheight,keepaspectratio]{#2}\end{center}}%
  {\CaptureManquante{#1}{\detokenize{#2}}}}
\newcommand{\CredixFigure}[1]{%
  \begin{center}\includegraphics[width=0.95\linewidth,height=0.8\textheight,keepaspectratio]{#1}\end{center}}
\usepackage{bookmark}
\IfFileExists{xurl.sty}{\usepackage{xurl}}{} % add URL line breaks if available
\urlstyle{same}
\hypersetup{
  pdftitle={Guide 2 : la méthode de conception},
  pdfauthor={SINGHE PENKA Hendrix Donavan},
  pdflang={fr},
  colorlinks=true,
  linkcolor={credixblue},
  filecolor={Maroon},
  citecolor={Blue},
  urlcolor={credixblue},
  pdfcreator={LaTeX via pandoc}}

\author{}
\date{}

\begin{document}

\begin{titlepage}
\thispagestyle{empty}
\centering
\vspace*{3.2cm}
{\color{credixgray}\large\textsc{Projet de fin d'études --- ENSPY, Université de Yaoundé I}\par}
\vspace{2.2cm}
{\Huge\bfseries\color{credixblue} Guide 2 : la méthode de conception\par}
\vspace{0.8cm}
{\Large Concevoir une application avec Claude, de A à Z\par}
\vspace{2cm}
{\color{credixblue}\rule{0.5\linewidth}{0.8pt}\par}
\vspace{1.2cm}
{\large SINGHE PENKA Hendrix Donavan\par}
{\color{credixgray}Dépôt GitHub : HendrixPenka2/credix\par}
\vfill
{\color{credixgray}\small Version du 19/09/2026\par}
\end{titlepage}


\renewcommand*\contentsname{Table des matières}
{
\setcounter{tocdepth}{2}
\tableofcontents
}
\section{Guide 2 : la méthode de travail avec Claude, de A à
Z}\label{guide-2--la-muxe9thode-de-travail-avec-claude-de-a-uxe0-z}

\textbf{Comment conduire un projet d\textquotesingle application avec
Claude, du besoin à l\textquotesingle application lancée. Illustrée par
la conception réelle de CREDIX.}

\begin{quote}
\textbf{Comment lire ce guide.} Les blocs intitulés \textbf{« Prompt »}
sont à copier dans Claude, après avoir remplacé ce qui est entre
crochets \texttt{{[}...{]}}. Les encadrés \textbf{« Dans CREDIX »}
montrent ce qui a été réellement fait, d\textquotesingle après les
documents du projet (instructions du projet, journaux de sessions et de
rédaction, rapports, documents de passation). Les prompts marqués
\textbf{« extrait réel »} sont tirés de ces fichiers ; les autres sont
des prompts types, construits sur le même modèle.

Ce guide suppose que vous avez lu le \textbf{guide 1} (abonnement, mode
Projet, Claude Code).
\end{quote}

\begin{center}\rule{0.5\linewidth}{0.5pt}\end{center}

\subsection{1. Le principe en une page}\label{1-le-principe-en-une-page}

Concevoir une application avec Claude ne consiste pas à lui dire «
fais-moi une application de crédit ». Cela donne un résultat fragile,
sans traçabilité. La méthode qui a servi pour CREDIX repose sur
\textbf{six principes}.

\begin{longtable}[]{@{}
  >{\raggedright\arraybackslash}p{(\linewidth - 2\tabcolsep) * \real{0.3471}}
  >{\raggedright\arraybackslash}p{(\linewidth - 2\tabcolsep) * \real{0.6229}}@{}}
\toprule\noalign{}
\begin{minipage}[b]{\linewidth}\raggedright
Principe
\end{minipage} & \begin{minipage}[b]{\linewidth}\raggedright
Ce que cela veut dire concrètement
\end{minipage} \\
\midrule\noalign{}
\endhead
\bottomrule\noalign{}
\endlastfoot
\textbf{1. On conçoit avant de coder} & Un cahier des charges, puis une
conception écrite, \textbf{avant} la première ligne de code. Règle du
projet professionnel : \emph{« ne jamais coder avant
d\textquotesingle avoir clarifié la conception »} \\
\textbf{2. Claude propose, l\textquotesingle humain décide} & Claude
explique ce qu\textquotesingle il va produire \textbf{avant} de le
produire. Aucun livrable final sans \textbf{accord explicite} \\
\textbf{3. Une seule tâche à la fois} & Le travail est découpé ; chaque
tâche suit les mêmes étapes \\
\textbf{4. Chaque décision est écrite} & Un registre consigne le choix,
les alternatives écartées, la justification et le \textbf{niveau de
preuve} \\
\textbf{5. Tout passe par des fichiers Markdown} & Le travail
d\textquotesingle une session devient un document, qui \textbf{alimente
la session suivante} \\
\textbf{6. On vérifie par l\textquotesingle exécution} & Chaque module
est testé ; les commandes de lancement sont rejouées sur une machine
neuve \\
\end{longtable}

\begin{quote}
\textbf{Pourquoi le Markdown (\texttt{.md}) ?} Claude écrit
naturellement en Markdown : du texte simple, avec titres, listes et
tableaux. Il est lisible par une personne \textbf{et} par un assistant
comme Claude Code, comparable ligne à ligne dans Git, et convertible
ensuite en PDF, en Word ou en présentation. \textbf{On garde toujours le
\texttt{.md} comme source}, même si le livrable final est un PDF (les
documents de ce dossier sont produits ainsi, avec LaTeX : voir le guide
4).
\end{quote}

\textbf{Le fil rouge de CREDIX, écrit dans le mémoire :} \emph{« chaque
bloc du pipeline sert un objectif précis, et on le prouve. »} La même
exigence vaut pour la méthode : \emph{chaque étape produit un document
vérifiable.}

\begin{center}\rule{0.5\linewidth}{0.5pt}\end{center}

\subsection{2. Vue d\textquotesingle ensemble : trois niveaux et neuf
étapes}\label{2-vue-densemble--trois-niveaux-et-neuf-uxe9tapes}

\subsubsection{2.1 Les trois niveaux de la
méthode}\label{21-les-trois-niveaux-de-la-muxe9thode}

\begin{longtable}[]{@{}
  >{\raggedright\arraybackslash}p{(\linewidth - 4\tabcolsep) * \real{0.2153}}
  >{\raggedright\arraybackslash}p{(\linewidth - 4\tabcolsep) * \real{0.5797}}
  >{\raggedright\arraybackslash}p{(\linewidth - 4\tabcolsep) * \real{0.1750}}@{}}
\toprule\noalign{}
\begin{minipage}[b]{\linewidth}\raggedright
Niveau
\end{minipage} & \begin{minipage}[b]{\linewidth}\raggedright
Ce que c\textquotesingle est
\end{minipage} & \begin{minipage}[b]{\linewidth}\raggedright
Où en parle-t-on
\end{minipage} \\
\midrule\noalign{}
\endhead
\bottomrule\noalign{}
\endlastfoot
\textbf{A. Les outils} & Un \textbf{Projet} claude.ai (instructions et
base de connaissances), \textbf{Claude Code} dans VS Code, des
\textbf{fichiers Markdown} & Guide 1 \\
\textbf{B. Le cycle de session} & Une façon
\textbf{d\textquotesingle ouvrir, de conduire et de fermer} chaque
séance de travail, pour ne jamais perdre le fil & Section 3 \\
\textbf{C. Les documents du projet} & Instructions, journal, registre de
décisions, pièges, TODO, documents de passation : la \textbf{mémoire} du
projet & Section 4 \\
\end{longtable}

\subsubsection{2.2 Les neuf étapes d\textquotesingle un
projet}\label{22-les-neuf-uxe9tapes-dun-projet}

\begin{longtable}[]{@{}
  >{\raggedright\arraybackslash}p{(\linewidth - 6\tabcolsep) * \real{0.1250}}
  >{\raggedright\arraybackslash}p{(\linewidth - 6\tabcolsep) * \real{0.3434}}
  >{\raggedright\arraybackslash}p{(\linewidth - 6\tabcolsep) * \real{0.2891}}
  >{\raggedright\arraybackslash}p{(\linewidth - 6\tabcolsep) * \real{0.2125}}@{}}
\toprule\noalign{}
\begin{minipage}[b]{\linewidth}\raggedright
Étape
\end{minipage} & \begin{minipage}[b]{\linewidth}\raggedright
Ce qu\textquotesingle on fait
\end{minipage} & \begin{minipage}[b]{\linewidth}\raggedright
Où
\end{minipage} & \begin{minipage}[b]{\linewidth}\raggedright
Livrable
\end{minipage} \\
\midrule\noalign{}
\endhead
\bottomrule\noalign{}
\endlastfoot
\textbf{0} & Préparer l\textquotesingle espace : Projet Claude,
instructions, journal & claude.ai & Projet configuré, journal \\
\textbf{1} & Passer du besoin au \textbf{cahier des charges} & claude.ai
(Projet) &
\texttt{cahier\_\allowbreak{}des\_\allowbreak{}charges.\allowbreak{}md} \\
\textbf{2} & Chercher dans la \textbf{littérature} & claude.ai
(recherche approfondie) + articles PDF & Fiches de lecture,
références \\
\textbf{3} & Faire les \textbf{choix méthodologiques} et les consigner &
claude.ai (Projet) & \textbf{Registre des décisions} \\
\textbf{4} & Choisir les \textbf{langages et technologies} & claude.ai
(Projet) & Section « choix techniques » \\
\textbf{5} & Rédiger le \textbf{document de conception} (UML, flux,
données) & claude.ai (Projet) & \texttt{conception.md} \\
\textbf{6} & Passer à \textbf{Claude Code} : mode \textbf{Plan}, TODO &
VS Code & Plan validé \\
\textbf{7} & \textbf{Prototyper module par module}, avec tests & VS Code
& Modules testés \\
\textbf{8} & \textbf{Livrer} : README de lancement, tests à la main & VS
Code + terminal & Projet lançable \\
\textbf{9} & \textbf{Modifier}, documenter, transmettre & VS Code (mode
Plan) & Documents de reprise \\
\end{longtable}

Chaque étape est conduite \textbf{à l\textquotesingle intérieur
d\textquotesingle une session}, selon le cycle de la section suivante.

\begin{center}\rule{0.5\linewidth}{0.5pt}\end{center}

\subsection{3. Le cycle de session : ouvrir, travailler,
fermer}\label{3-le-cycle-de-session--ouvrir-travailler-fermer}

Un projet long se déroule sur \textbf{de nombreuses sessions}. Claude ne
se souvient pas d\textquotesingle une session à l\textquotesingle autre
: c\textquotesingle est \textbf{vous} qui organisez la continuité, avec
des documents. Toutes les sessions de CREDIX suivent le même cycle.

\subsubsection{3.1 Ouvrir la session : le prompt de
démarrage}\label{31-ouvrir-la-session--le-prompt-de-duxe9marrage}

\textbf{Principe.} On ne pose pas de question tout de suite. Le premier
message impose \textbf{ce que Claude doit lire, dans quel ordre}, et lui
interdit de commencer avant votre accord.

\textbf{Prompt de démarrage (extrait réel, phase de rédaction du
chapitre 2 de CREDIX)}

\begin{Verbatim}[breaklines=true,breakanywhere=true,fontsize=\footnotesize,frame=leftline,framerule=1.6pt,rulecolor=\color{credixblue},framesep=3mm,xleftmargin=2mm,xrightmargin=1mm]
Voici la suite du travail sur le mémoire CREDIX, chapitre 2. Avant de répondre
quoi que ce soit, lis dans cet ordre, entièrement :

1. HANDOVER_SESSION_SUIVANTE.md : à jour, fin de la session précédente.
2. JOURNAL_REDACTION_CREDIX_CH2_A_JOUR.md : lis en entier la section « Pièges
   méthodologiques déjà identifiés » avant toute chose : ce ne sont pas des
   recommandations générales, ce sont des erreurs déjà commises dans ce projet.
3. chap02_en_cours.tex : le chapitre en cours.
4. TODO_CHAPITRE2_CREDIX_MAJ.md : le plan détaillé à jour. Le journal fait foi
   en cas de désaccord entre les deux fichiers.
5. Dans le project knowledge : [les rapports à lire en entier avant d'écrire].
6. Le notebook joint : nécessaire pour vérifier tout chiffre cité.

Avant d'ouvrir la section suivante, j'ai des remarques groupées sur ma relecture :
je te les donne d'abord, on les traite, puis on passe à la suite. Ne commence pas
la rédaction avant que je te dise que c'est bon.
\end{Verbatim}

\textbf{Ce que ce prompt impose, et pourquoi :}

\begin{longtable}[]{@{}
  >{\raggedright\arraybackslash}p{(\linewidth - 2\tabcolsep) * \real{0.4311}}
  >{\raggedright\arraybackslash}p{(\linewidth - 2\tabcolsep) * \real{0.5389}}@{}}
\toprule\noalign{}
\begin{minipage}[b]{\linewidth}\raggedright
Élément
\end{minipage} & \begin{minipage}[b]{\linewidth}\raggedright
Effet
\end{minipage} \\
\midrule\noalign{}
\endhead
\bottomrule\noalign{}
\endlastfoot
\textbf{Ordre de lecture explicite} & Claude part du même état que vous,
sans oublier les documents essentiels \\
\textbf{Le journal « fait foi »} en cas de désaccord & Il y a
\textbf{une} source de vérité \\
\textbf{Lire d\textquotesingle abord les pièges} & On évite de
\textbf{refaire} des erreurs déjà commises \\
\textbf{« Ne commence pas avant que je te dise que c\textquotesingle est
bon »} & Vous gardez la main : Claude propose, vous décidez \\
\end{longtable}

\textbf{Variante pour un projet de programmation (extrait réel, projet
professionnel)}

\begin{Verbatim}[breaklines=true,breakanywhere=true,fontsize=\footnotesize,frame=leftline,framerule=1.6pt,rulecolor=\color{credixblue},framesep=3mm,xleftmargin=2mm,xrightmargin=1mm]
Tu es mon pair-programmer principal. Règles absolues : tu travailles une seule
tâche à la fois ; tu respectes l'architecture existante et les conventions ;
tu ne codes jamais avant d'avoir clarifié la conception.
Méthode obligatoire en 5 étapes : 1) analyse de l'existant, 2) conception
fonctionnelle, 3) conception technique, 4) mise en place, 5) tests.
Je vais te fournir le contexte (fichiers, TODO, éventuellement un rapport de
passation). Attends mes instructions.
\end{Verbatim}

Et le message de démarrage de ce même kit : \emph{« Ta première mission
AVANT toute action : 1. lire CONTEXT.md pour comprendre la stack et les
fichiers clés ; 2. lire TODO.md et identifier les tâches non cochées ;
3. me dire quelle est la prochaine tâche logique ; 4. attendre ma
validation avant de commencer. \textbf{Ne propose aucun code pour
l\textquotesingle instant.} »}

\subsubsection{3.2 Conduire la session : une tâche à la fois, en cinq
étapes}\label{32-conduire-la-session--une-tuxe2che-uxe0-la-fois-en-cinq-uxe9tapes}

Pendant la session, on respecte les règles inscrites dans les
instructions du projet (voir section 4.1) :

\begin{itemize}
\tightlist
\item
  \textbf{Claude explique chaque initiative avant de produire quoi que
  ce soit.}
\item
  \textbf{On travaille sur une seule tâche}, définie au début de la
  session, et on ne touche pas au reste. Exemple réel : \emph{« Cette
  session a un seul objectif : exécuter le bloc A de la TODO. Ne pas
  toucher au backend. Ne pas travailler sur le mémoire. Uniquement le
  notebook. »}
\item
  Pour une tâche technique, on suit les \textbf{cinq étapes} :
  \textbf{analyse de l\textquotesingle existant → conception
  fonctionnelle → conception technique → mise en place → tests}.
\item
  Pour une tâche de rédaction, on procède \textbf{sous-section par
  sous-section} : décider ce qui entre, ce qui sort, ce qui part dans un
  autre chapitre ; proposer un sous-plan ; discuter le contenu ; rédiger
  ; faire une vérification technique courte ; réserver
  l\textquotesingle{}\textbf{audit critique complet} aux sous-sections à
  risque, identifiées ensemble.
\item
  \textbf{Claude signale quand le contexte devient chargé} et propose de
  changer de session. Si vous acceptez, il prépare la documentation de
  clôture.
\end{itemize}

\subsubsection{3.3 Fermer la session : le rapport de
passation}\label{33-fermer-la-session--le-rapport-de-passation}

\textbf{Prompt de fin de session (extrait réel, projet professionnel)}

\begin{Verbatim}[breaklines=true,breakanywhere=true,fontsize=\footnotesize,frame=leftline,framerule=1.6pt,rulecolor=\color{credixblue},framesep=3mm,xleftmargin=2mm,xrightmargin=1mm]
Nous allons bientôt changer de session. Génère-moi un « Handoff Report »
(résumé de passation) très concis. Il doit contenir :
1. Ce qui a été complété avec succès aujourd'hui.
2. Ce qui est actuellement en cours (avec le fichier exact et la ligne si possible).
3. La prochaine étape immédiate selon notre méthode en 5 étapes.
4. Les éventuels points d'attention ou bugs laissés en suspens.
Ne génère pas de code, juste ce rapport en markdown.
\end{Verbatim}

\textbf{La règle de clôture des instructions du projet de CREDIX :}
\emph{les livrables de clôture ne sont produits qu\textquotesingle après
l\textquotesingle accord explicite de l\textquotesingle auteur, et
Claude explique ce qu\textquotesingle il va produire avant de le
produire.} Ce sont :

\begin{enumerate}
\def\labelenumi{\arabic{enumi}.}
\tightlist
\item
  les \textbf{instructions du projet} mises à jour, si elles ont changé
  ;
\item
  une \textbf{nouvelle entrée du journal de bord}, datée, en
  \textbf{ajout seulement} (on ne réécrit jamais le passé) ;
\item
  le \textbf{plan} et la \textbf{TODO} mis à jour ;
\item
  le \textbf{document de passation}, qui devient
  l\textquotesingle entrée de la session suivante.
\end{enumerate}

Ces documents sont ensuite \textbf{déposés dans le projet Claude} (ou
dans le dossier du projet de code) : c\textquotesingle est ce qui rend
la session suivante possible.

\begin{quote}
\textbf{Dans CREDIX.} Le résultat est visible dans le dépôt : des
documents de passation (\texttt{HANDOVER\_...}), un journal de bord, des
rapports de session par sujet (par exemple le rapport de la séance sur
le flux de contrôle des profils atypiques), une « copie miroir » des
notes que Claude conserve d\textquotesingle une session à
l\textquotesingle autre, et des prompts de démarrage propres à chaque
type de séance (notebook, rédaction du mémoire, codage).
\end{quote}

\begin{center}\rule{0.5\linewidth}{0.5pt}\end{center}

\subsection{4. Les documents qui font tenir le
projet}\label{4-les-documents-qui-font-tenir-le-projet}

\subsubsection{4.1 Les instructions du projet (six
sections)}\label{41-les-instructions-du-projet-six-sections}

C\textquotesingle est le document le plus important : il est chargé
\textbf{une fois} dans le Projet Claude et s\textquotesingle applique à
\textbf{toutes} les conversations. Voici la structure réelle de celui de
CREDIX.

\begin{longtable}[]{@{}
  >{\raggedright\arraybackslash}p{(\linewidth - 4\tabcolsep) * \real{0.2125}}
  >{\raggedright\arraybackslash}p{(\linewidth - 4\tabcolsep) * \real{0.4494}}
  >{\raggedright\arraybackslash}p{(\linewidth - 4\tabcolsep) * \real{0.3081}}@{}}
\toprule\noalign{}
\begin{minipage}[b]{\linewidth}\raggedright
Section
\end{minipage} & \begin{minipage}[b]{\linewidth}\raggedright
Contenu
\end{minipage} & \begin{minipage}[b]{\linewidth}\raggedright
Pourquoi
\end{minipage} \\
\midrule\noalign{}
\endhead
\bottomrule\noalign{}
\endlastfoot
\textbf{1. Contexte} & Qui, quoi, pour qui, dates clés, encadrants,
sujet, état d\textquotesingle avancement & Claude sait dans quel cadre
il travaille \\
\textbf{2. Ressources permanentes} & La liste \textbf{classée} des
documents du projet, avec pour chacun \textbf{« consulter pour\ldots{}
»} : modèles de référence, documents techniques, normes de rédaction,
norme bibliographique, modèle LaTeX, version courante du mémoire &
Claude sait \textbf{où chercher} selon la tâche \\
\textbf{3. Décisions techniques figées} & Treize décisions prises, à ne
remettre en question que sur contradiction logique majeure (modèle
final, ordre des étapes du pipeline, traitement des manquants, découpage
des données, métriques principales, pile technique\ldots) & Évite de
rediscuter sans cesse les mêmes choix \\
\textbf{4. Consignes permanentes} & Langue, niveau attendu, style,
règles anti-plagiat, outil pour les diagrammes, cohérence code/rapport ;
\textbf{« à faire systématiquement »} et \textbf{« ce que Claude ne doit
pas faire »} & Uniformise toutes les réponses \\
\textbf{5. Workflow de collaboration} & Ce qui se passe \textbf{au
début, pendant et à la fin} de chaque session (voir section 3) & Le
cycle de session est écrit noir sur blanc \\
\textbf{6. État courant} & Où en est le projet, mis à jour à chaque
session & Claude connaît la situation \\
\end{longtable}

\textbf{Extrait réel de la section 4 (consignes permanentes) :}

\begin{Verbatim}[breaklines=true,breakanywhere=true,fontsize=\footnotesize,frame=leftline,framerule=1.6pt,rulecolor=\color{credixblue},framesep=3mm,xleftmargin=2mm,xrightmargin=1mm]
À faire systématiquement :
- Lire le journal de bord en début de chaque session
- Consulter [le document de méthodologie] avant tout chapitre technique
- Consulter [le rapport de pipeline] pour tous les résultats numériques
- Signaler immédiatement tout doublon détecté dans la base de connaissances
- Demander l'OK explicite de l'auteur avant de produire les livrables finaux

Ce que Claude ne doit pas faire :
- Rédiger sans avoir consulté les ressources pertinentes
- Produire des livrables sans OK explicite de l'auteur
\end{Verbatim}

\textbf{Extrait réel de la section 5 (workflow) :}

\begin{Verbatim}[breaklines=true,breakanywhere=true,fontsize=\footnotesize,frame=leftline,framerule=1.6pt,rulecolor=\color{credixblue},framesep=3mm,xleftmargin=2mm,xrightmargin=1mm]
En début de session : Claude lit le journal ; propose un objectif de session ;
l'auteur confirme ou ajuste.
Pendant la session : on itère ; Claude explique chaque initiative avant de produire.
En fin de session : Claude produit les livrables de clôture UNIQUEMENT après OK
explicite : instructions à jour, journal (nouvelle entrée datée, ajout seulement),
plans à jour.
Règle fondamentale : aucun livrable final n'est produit sans OK explicite.
\end{Verbatim}

\subsubsection{4.2 Les autres documents, en un
tableau}\label{42-les-autres-documents-en-un-tableau}

\begin{longtable}[]{@{}
  >{\raggedright\arraybackslash}p{(\linewidth - 6\tabcolsep) * \real{0.2500}}
  >{\raggedright\arraybackslash}p{(\linewidth - 6\tabcolsep) * \real{0.2912}}
  >{\raggedright\arraybackslash}p{(\linewidth - 6\tabcolsep) * \real{0.2038}}
  >{\raggedright\arraybackslash}p{(\linewidth - 6\tabcolsep) * \real{0.2250}}@{}}
\toprule\noalign{}
\begin{minipage}[b]{\linewidth}\raggedright
Document
\end{minipage} & \begin{minipage}[b]{\linewidth}\raggedright
Rôle
\end{minipage} & \begin{minipage}[b]{\linewidth}\raggedright
Règle d\textquotesingle or
\end{minipage} & \begin{minipage}[b]{\linewidth}\raggedright
Exemple dans CREDIX
\end{minipage} \\
\midrule\noalign{}
\endhead
\bottomrule\noalign{}
\endlastfoot
\textbf{Journal de bord} & Une entrée par session : \emph{réalisé,
décisions prises, en attente} & \textbf{Ajout seulement}, jamais de
réécriture du passé &
\texttt{journal\_\allowbreak{}sessions\_\allowbreak{}memoire\_\allowbreak{}hendrix.\allowbreak{}txt} \\
\textbf{Journal de rédaction} & État consolidé d\textquotesingle un
chapitre : plan, \textbf{règles de rédaction actées}, \textbf{décisions
verrouillées}, \textbf{pièges}, points ouverts & \textbf{Fait foi} en
cas de désaccord avec la TODO &
\texttt{JOURNAL\_\allowbreak{}REDACTION\_\allowbreak{}CREDIX\_\allowbreak{}CH2\_\allowbreak{}A\_\allowbreak{}JOUR.\allowbreak{}md} \\
\textbf{Registre des décisions} & Chaque choix : options, justification,
niveau de preuve & Toute décision importante y figure &
\texttt{RAPPORT\_\allowbreak{}DECISIONS\_\allowbreak{}METHODOLOGIQUES\_\allowbreak{}CREDIX.\allowbreak{}md} \\
\textbf{Rapport de résultats} & Résultats validés, consolidés &
\textbf{Aucun chiffre inventé} : il vient d\textquotesingle un calcul &
\texttt{RAPPORT\_\allowbreak{}RESULTATS\_\allowbreak{}CONCLUSIONS\_\allowbreak{}CREDIX.\allowbreak{}md} \\
\textbf{Pièges méthodologiques} & La liste des \textbf{erreurs déjà
commises} & À relire \textbf{avant} de reprendre le travail & Section
4bis du journal de rédaction \\
\textbf{TODO} & Le plan détaillé, tâche par tâche & Une tâche = une case
à cocher & \texttt{TODO\_...md} \\
\textbf{Document de passation} & Ce qui est fait, en cours, à faire,
règles absolues & Écrit à la fin de \textbf{chaque} session ou phase &
\texttt{HANDOVER\_...md} \\
\textbf{État miroir} & Copie des notes que Claude garde
d\textquotesingle une session à l\textquotesingle autre & Doit rester
identique à la mémoire interne &
\texttt{ETAT\_\allowbreak{}MEMOIRE\_\allowbreak{}CLAUDE\_\allowbreak{}CREDIX.\allowbreak{}md} \\
\textbf{Instantané de contexte} & Un fichier qui résume \textbf{tout le
code} d\textquotesingle un projet, fabriqué par un script & À régénérer
avant d\textquotesingle en parler à Claude & \texttt{scanner.sh} produit
\texttt{contexte\_\allowbreak{}fastapi.\allowbreak{}txt} \\
\end{longtable}

\subsubsection{4.3 Deux exemples réels qui montrent la valeur de ces
documents}\label{43-deux-exemples-ruxe9els-qui-montrent-la-valeur-de-ces-documents}

\textbf{Des règles de rédaction « actées »} (extraits du journal de
rédaction, sur trente-huit règles) :

\begin{itemize}
\tightlist
\item
  \emph{« Un paragraphe = une fonction argumentative. »}
\item
  \emph{« Aucun composant sans fonction énonçable en une phrase. »}
\item
  \emph{« Séparation stricte conception / résultats : jamais "meilleure"
  avant les résultats. »}
\item
  \emph{« \textbf{Grille des quatre questions pour chaque bloc
  technique} : ① le problème ② la solution ③ pourquoi celle-là ④ comment
  elle est validée. »}
\item
  \emph{« Terminologie : un seul terme, "dispositif de décision". Jamais
  "dispositif de contrôle". »}
\end{itemize}

\textbf{Des pièges consignés, avec leur histoire} (extraits de la
section « Pièges méthodologiques ») :

\begin{itemize}
\tightlist
\item
  \textbf{Vérifier tout chiffre reçu d\textquotesingle une source
  externe, même quand le raisonnement est juste.} Deux documents
  externes utilisaient une variable comme exemple, avec un indicateur
  d\textquotesingle une valeur de 0,13. Vérifié contre
  l\textquotesingle exécution réelle : la valeur réelle était de
  \textbf{0,0199}, dans une autre catégorie. Le raisonnement général
  restait correct et a été conservé ; \textbf{seul le chiffre était
  faux}.
\item
  \textbf{Ne jamais présenter le système autrement qu\textquotesingle il
  n\textquotesingle est}, « même sur un détail », pour paraître plus
  conforme à la littérature.
\item
  \textbf{Ne jamais dire « on a gardé le moins bon »} : une décision se
  justifie par des critères, pas par un aveu.
\end{itemize}

\subsubsection{4.4 Comment les créer avec
Claude}\label{44-comment-les-cruxe9er-avec-claude}

\textbf{Prompt : rédiger le journal ou le document de passation}

\begin{Verbatim}[breaklines=true,breakanywhere=true,fontsize=\footnotesize,frame=leftline,framerule=1.6pt,rulecolor=\color{credixblue},framesep=3mm,xleftmargin=2mm,xrightmargin=1mm]
Fais le document de reprise de cette session en Markdown : 1) l'objectif de la
session, 2) ce qui a été fait, 3) les décisions prises et leur justification,
4) ce qui a été écarté et pourquoi, 5) les fichiers créés ou modifiés,
6) ce qui reste à faire dans l'ordre, 7) les pièges à éviter. Sois factuel :
aucun chiffre que nous n'avons pas obtenu.
\end{Verbatim}

Les modèles prêts à remplir sont en section 16.

\begin{center}\rule{0.5\linewidth}{0.5pt}\end{center}

\subsection{5. Étape 0 : préparer l\textquotesingle espace de
travail}\label{5-uxe9tape-0--pruxe9parer-lespace-de-travail}

\subsubsection{5.1 Ce qu\textquotesingle il faut mettre en
place}\label{51-ce-quil-faut-mettre-en-place}

\begin{enumerate}
\def\labelenumi{\arabic{enumi}.}
\tightlist
\item
  \textbf{Un Projet Claude} (guide 1, section 6), avec ses
  \textbf{instructions en six sections} (section 4.1 de ce guide).
\item
  \textbf{Un journal de bord}, en ajout seulement.
\item
  \textbf{Un dossier de projet} sur l\textquotesingle ordinateur, sous
  \textbf{Git}, dès le début.
\end{enumerate}

\subsubsection{5.2 La toute première
session}\label{52-la-toute-premiuxe8re-session}

\begin{quote}
\textbf{Dans CREDIX (27 mai 2026).} La première session a servi à
\textbf{initialiser le projet Claude} : inventaire des ressources de la
base de connaissances (cahier des charges v2.0, conception v2.1,
méthodologie complète v2, rapport de pipeline, modèles de mémoire de
référence), \textbf{diagnostic} de l\textquotesingle existant (pipeline
documenté à environ 75 \%, backend à 0 \%, frontend à environ 15 \%),
rédaction des instructions du projet, d\textquotesingle un plan du
mémoire, d\textquotesingle un \textbf{plan de travail daté} et du
\textbf{journal de bord}.
\end{quote}

\textbf{Prompt 0 : faire l\textquotesingle inventaire du projet}

\begin{Verbatim}[breaklines=true,breakanywhere=true,fontsize=\footnotesize,frame=leftline,framerule=1.6pt,rulecolor=\color{credixblue},framesep=3mm,xleftmargin=2mm,xrightmargin=1mm]
Fais l'inventaire de tous les documents de la base de connaissances de ce projet.
Pour chacun : son titre, son rôle, sa date, et s'il est à jour ou périmé.
Puis fais un diagnostic honnête : ce qui est prêt, ce qui manque, ce qui se
contredit. Termine par les trois prochaines actions que tu recommandes.
\end{Verbatim}

\CredixCapture{Base de connaissances d'un projet avec les documents chargés}{../images/guides/02-01_base_connaissances.png}

\emph{Figure 1. La base de connaissances du projet.}

\begin{center}\rule{0.5\linewidth}{0.5pt}\end{center}

\subsection{6. Étape 1 : du besoin au cahier des
charges}\label{6-uxe9tape-1--du-besoin-au-cahier-des-charges}

\subsubsection{6.1 Laisser Claude vous
interroger}\label{61-laisser-claude-vous-interroger}

Le piège du débutant est de rédiger seul un long descriptif. Il vaut
mieux \textbf{demander à Claude de vous poser des questions}, une à la
fois.

\textbf{Prompt 1 : cadrer le besoin par des questions}

\begin{Verbatim}[breaklines=true,breakanywhere=true,fontsize=\footnotesize,frame=leftline,framerule=1.6pt,rulecolor=\color{credixblue},framesep=3mm,xleftmargin=2mm,xrightmargin=1mm]
Je veux concevoir [décrire l'application en une phrase], pour [qui].
Avant de rédiger quoi que ce soit, pose-moi une question à la fois pour
clarifier : le problème à résoudre, les utilisateurs, ce que l'application
doit faire, ce qu'elle ne doit pas faire, les contraintes (temps, budget,
données, règles légales). Attends ma réponse avant la question suivante.
Quand tu estimes avoir assez d'éléments, dis-le-moi.
\end{Verbatim}

\subsubsection{6.2 Rédiger et faire critiquer le cahier des
charges}\label{62-ruxe9diger-et-faire-critiquer-le-cahier-des-charges}

\textbf{Prompt 2 : rédiger le cahier des charges}

\begin{Verbatim}[breaklines=true,breakanywhere=true,fontsize=\footnotesize,frame=leftline,framerule=1.6pt,rulecolor=\color{credixblue},framesep=3mm,xleftmargin=2mm,xrightmargin=1mm]
À partir de notre échange, rédige le cahier des charges en Markdown avec :
1. Contexte et objectifs
2. Acteurs et leurs rôles
3. Périmètre (inclus) et hors périmètre (exclu)
4. Exigences fonctionnelles, numérotées EF-1, EF-2...
5. Exigences non fonctionnelles (performance, sécurité, traçabilité...), ENF-1...
6. Contraintes et hypothèses
7. Critères d'acceptation : pour chaque exigence, comment vérifier qu'elle est atteinte
Exigences courtes, vérifiables, sans ambiguïté.
\end{Verbatim}

\textbf{Prompt 3 : faire critiquer le cahier des charges}

\begin{Verbatim}[breaklines=true,breakanywhere=true,fontsize=\footnotesize,frame=leftline,framerule=1.6pt,rulecolor=\color{credixblue},framesep=3mm,xleftmargin=2mm,xrightmargin=1mm]
Relis ce cahier des charges comme un examinateur sévère. Liste : les ambiguïtés,
les contradictions, les exigences impossibles à vérifier, ce qui manque.
Classe tes remarques par gravité. Ne corrige rien : propose seulement.
\end{Verbatim}

\begin{quote}
\textbf{Dans CREDIX.} Les exigences du mémoire sont numérotées (EM-1 à
EM-9). Un cahier des charges séparé a été produit pour le frontend
(\texttt{cdc\_credix.md} : produit, trois rôles, lexique métier, plan de
navigation de 36 routes) avant de dessiner les écrans.
\end{quote}

\begin{center}\rule{0.5\linewidth}{0.5pt}\end{center}

\subsection{7. Étape 2 : la recherche dans la
littérature}\label{7-uxe9tape-2--la-recherche-dans-la-littuxe9rature}

Un choix technique important s\textquotesingle appuie sur la
\textbf{littérature}. Claude peut faire une grande partie de ce travail,
à condition de \textbf{ne rien croire sans vérifier}. On combine trois
outils :

\begin{enumerate}
\def\labelenumi{\arabic{enumi}.}
\tightlist
\item
  la \textbf{recherche approfondie} de claude.ai (guide 1, section 5.2),
  qui produit un rapport \textbf{avec citations} ;
\item
  les \textbf{articles PDF} chargés dans le Projet : Claude répond
  \textbf{à partir d\textquotesingle eux} ;
\item
  \textbf{votre vérification} : ouvrir les références clés.
\end{enumerate}

\textbf{Prompt 4 : balayer la littérature}

\begin{Verbatim}[breaklines=true,breakanywhere=true,fontsize=\footnotesize,frame=leftline,framerule=1.6pt,rulecolor=\color{credixblue},framesep=3mm,xleftmargin=2mm,xrightmargin=1mm]
Fais une recherche approfondie sur : [question précise, par exemple
« méthodes de sélection de variables pour un scoring de crédit explicable »].
Je veux : les approches principales, leurs avantages et limites, les références
les plus citées (auteur, année, titre). Cite tes sources. Indique ce qui est un
consensus et ce qui est discuté. Termine par les références dont tu n'es pas
sûr, marquées [à vérifier].
\end{Verbatim}

\textbf{Prompt 5 : interroger un article chargé dans le Projet}

\begin{Verbatim}[breaklines=true,breakanywhere=true,fontsize=\footnotesize,frame=leftline,framerule=1.6pt,rulecolor=\color{credixblue},framesep=3mm,xleftmargin=2mm,xrightmargin=1mm]
Dans l'article [titre ou auteur] chargé dans ce projet :
1. Quelle est la méthode proposée, en termes simples ?
2. Sur quelles données est-elle évaluée, et avec quels résultats ?
3. Quelles sont ses limites, d'après les auteurs et d'après toi ?
4. Que peut-on en retenir pour mon choix de [sujet] ?
Cite les passages sur lesquels tu t'appuies.
\end{Verbatim}

\textbf{Les garde-fous :} vérifier chaque référence importante ; faire
marquer les références incertaines par \texttt{{[}à\ vérifier{]}} ;
\textbf{distinguer} ce que dit la littérature, ce qu\textquotesingle ont
montré vos essais et ce qui est un choix assumé (échelle de preuve,
section 8).

\begin{quote}
\textbf{Dans CREDIX.} Une session de recherche approfondie a été
conservée sous forme de PDF. Pour le seuil de 50 \% de valeurs
manquantes, la recherche a établi qu\textquotesingle{}\textbf{il
n\textquotesingle existe pas de règle universelle} : le package de
référence \texttt{scorecard} (langage R) fixe 95 \% par défaut, et la
littérature utilise des seuils très variables. Conclusion consignée :
\emph{ne pas attribuer ce seuil à un auteur} ; le présenter comme un
\textbf{choix assumé}. Des références comme Chow (1970) ou El-Yaniv et
Wiener (2010) sont restées marquées « à vérifier » tant
qu\textquotesingle elles n\textquotesingle étaient pas contrôlées.
\end{quote}

\CredixCapture{Rapport de recherche approfondie avec ses citations}{../images/guides/02-02_recherche_approfondie.png}

\emph{Figure 2. Un résultat de recherche approfondie avec citations.}

\begin{center}\rule{0.5\linewidth}{0.5pt}\end{center}

\subsection{8. Étape 3 : les choix méthodologiques, exemple du
WoE}\label{8-uxe9tape-3--les-choix-muxe9thodologiques-exemple-du-woe}

C\textquotesingle est le cœur de la méthode. Chaque choix suit la même
séquence : \textbf{question, options, critères, décision, preuve,
consignation.}

\subsubsection{8.1 L\textquotesingle échelle de preuve à trois
niveaux}\label{81-luxe9chelle-de-preuve-uxe0-trois-niveaux}

\begin{longtable}[]{@{}
  >{\raggedright\arraybackslash}p{(\linewidth - 4\tabcolsep) * \real{0.2000}}
  >{\raggedright\arraybackslash}p{(\linewidth - 4\tabcolsep) * \real{0.4244}}
  >{\raggedright\arraybackslash}p{(\linewidth - 4\tabcolsep) * \real{0.3456}}@{}}
\toprule\noalign{}
\begin{minipage}[b]{\linewidth}\raggedright
Niveau
\end{minipage} & \begin{minipage}[b]{\linewidth}\raggedright
Signification
\end{minipage} & \begin{minipage}[b]{\linewidth}\raggedright
Exemple
\end{minipage} \\
\midrule\noalign{}
\endhead
\bottomrule\noalign{}
\endlastfoot
\textbf{I : démontré} & Établi par une expérience sur nos données
(mesure, ablation, partition) & « L\textquotesingle AUC augmente quand
la couverture ρc augmente » \\
\textbf{II : ancré dans la littérature} & Fait connu, cité, non refait &
« Le WoE est standard en scoring de crédit (Siddiqi) » \\
\textbf{III : hypothèse assumée} & Choix de conception justifié comme
tel & « Seuil de refus à 30 \% de probabilité de défaut » \\
\end{longtable}

Et la \textbf{grille des quatre questions} pour chaque bloc technique :
\textbf{① le problème, ② la solution, ③ pourquoi celle-là, ④ comment
elle est validée.}

\subsubsection{8.2 L\textquotesingle exemple complet : pourquoi le WoE
?}\label{82-lexemple-complet--pourquoi-le-woe-}

La séquence ci-dessous est \textbf{reconstituée à partir des décisions
réellement consignées} dans les documents du projet.

\textbf{A. Poser la question, avec les contraintes.}

\begin{Verbatim}[breaklines=true,breakanywhere=true,fontsize=\footnotesize,frame=leftline,framerule=1.6pt,rulecolor=\color{credixblue},framesep=3mm,xleftmargin=2mm,xrightmargin=1mm]
Mon pipeline de scoring doit rester explicable variable par variable (contrainte
réglementaire) et gérer des valeurs manquantes nombreuses. Comment encoder mes
variables numériques et catégorielles avant le modèle ? Compare au moins :
standardisation + encodage one-hot, encodage par la cible, et le WoE.
Critères : explicabilité, traitement des manquants, stabilité, risque de fuite
de données. Termine par une recommandation et ses risques.
\end{Verbatim}

\textbf{B. Claude compare.} Le WoE (\emph{Weight of Evidence}) découpe
chaque variable en tranches et remplace chaque valeur par le logarithme
du rapport entre bons et mauvais payeurs de sa tranche.

\begin{longtable}[]{@{}
  >{\raggedright\arraybackslash}p{(\linewidth - 6\tabcolsep) * \real{0.2294}}
  >{\raggedright\arraybackslash}p{(\linewidth - 6\tabcolsep) * \real{0.3478}}
  >{\raggedright\arraybackslash}p{(\linewidth - 6\tabcolsep) * \real{0.2375}}
  >{\raggedright\arraybackslash}p{(\linewidth - 6\tabcolsep) * \real{0.1554}}@{}}
\toprule\noalign{}
\begin{minipage}[b]{\linewidth}\raggedright
Critère
\end{minipage} & \begin{minipage}[b]{\linewidth}\raggedright
WoE
\end{minipage} & \begin{minipage}[b]{\linewidth}\raggedright
One-hot + standardisation
\end{minipage} & \begin{minipage}[b]{\linewidth}\raggedright
Encodage par la cible
\end{minipage} \\
\midrule\noalign{}
\endhead
\bottomrule\noalign{}
\endlastfoot
Explicabilité & \textbf{Très bonne} : chaque tranche a un poids lisible
& Moyenne & Faible \\
Valeurs manquantes & \textbf{Tranche « Manquant » dédiée} & Imputation
arbitraire & À gérer à part \\
Stabilité (tranches monotones) & \textbf{Contrainte possible} & Non &
Non \\
Risque de fuite & Réel : à apprendre \textbf{sur le train seul} & Faible
& Élevé \\
\end{longtable}

\textbf{C. L\textquotesingle avocat du diable.}

\begin{Verbatim}[breaklines=true,breakanywhere=true,fontsize=\footnotesize,frame=leftline,framerule=1.6pt,rulecolor=\color{credixblue},framesep=3mm,xleftmargin=2mm,xrightmargin=1mm]
Donne-moi les cinq meilleures objections à ce choix. Pour chacune, dis si elle
est décisive ou si on peut la traiter.
\end{Verbatim}

\textbf{D. La décision, écrite dans le registre} (faits tirés du
récapitulatif de méthodologie du projet) :

\begin{longtable}[]{@{}
  >{\raggedright\arraybackslash}p{(\linewidth - 2\tabcolsep) * \real{0.2399}}
  >{\raggedright\arraybackslash}p{(\linewidth - 2\tabcolsep) * \real{0.7301}}@{}}
\toprule\noalign{}
\begin{minipage}[b]{\linewidth}\raggedright
Champ
\end{minipage} & \begin{minipage}[b]{\linewidth}\raggedright
Contenu
\end{minipage} \\
\midrule\noalign{}
\endhead
\bottomrule\noalign{}
\endlastfoot
\textbf{Décision} & Encodage \textbf{WoE} réalisé par la bibliothèque
\texttt{optbinning} (classe \texttt{OptimalBinning}) \\
\textbf{Pourquoi} & Explicabilité par variable exigée par la
réglementation ; gestion native des manquants ; tranches monotones \\
\textbf{Mécanisme} & Pré-découpage fin par arbre de décision, puis un
solveur de programmation par contraintes fusionne les tranches voisines
pour maximiser l\textquotesingle IV, sous contraintes de taille
minimale, de nombre maximal de tranches et de \textbf{monotonie du
WoE} \\
\textbf{Valeurs manquantes} & Jamais imputées : regroupées dans une
tranche \textbf{« Manquant »} dont le WoE est calculé sur la vraie
proportion de bons et de mauvais payeurs \\
\textbf{Garde-fou anti-fuite} & Découpage appris \textbf{sur le jeu
d\textquotesingle entraînement uniquement}, puis appliqué à la
validation et au test \\
\textbf{Filtre associé} & Variables retenues si \textbf{IV ≥ 0,02}
(seuil classique) \\
\textbf{Choix assumé} & Une variable avec \textbf{plus de 50 \% de
manquants} est exclue, sauf si son IV atteint \textbf{0,10} : plus
conservateur que le défaut de certains outils (95 \%) \\
\textbf{Niveau de preuve} & \textbf{II} pour le WoE et le seuil
d\textquotesingle IV ; \textbf{III} pour le seuil de 50 \% (à présenter
comme tel) \\
\textbf{Alternatives écartées} & Standardisation + one-hot ; encodage
par la cible \\
\textbf{Risques restants} & Perte d\textquotesingle information par
découpage en tranches \\
\end{longtable}

\textbf{E. Vérifier par une expérience quand c\textquotesingle est
possible.} Un contrôle cellule par cellule du notebook a révélé un
\textbf{défaut réel} : pour recalculer l\textquotesingle IV des colonnes
exclues, le code associait la valeur d\textquotesingle une colonne
(ordre brut du fichier) à la cible du découpage temporel (ordre trié).
Les deux ordres différaient : l\textquotesingle IV de ces colonnes ne
mesurait rien de fiable (valeurs quasi nulles pour une quarantaine de
colonnes). La correction (fusion sur l\textquotesingle identifiant
client) a fait passer 20 colonnes dans la catégorie « couverture faible
», \textbf{prévue mais jamais déclenchée}. Sans cette vérification, un
résultat faux serait entré dans le mémoire.

\subsubsection{8.3 Les mêmes réflexes pour tout autre
choix}\label{83-les-muxeames-ruxe9flexes-pour-tout-autre-choix}

La même séquence a servi pour : LightGBM contre XGBoost,
l\textquotesingle autoencodeur pour le contrôle des profils atypiques,
la calibration isotonique, les seuils de décision,
l\textquotesingle indice de couverture ρc.

\textbf{Prompt 6 : produire une fiche de décision}

\begin{Verbatim}[breaklines=true,breakanywhere=true,fontsize=\footnotesize,frame=leftline,framerule=1.6pt,rulecolor=\color{credixblue},framesep=3mm,xleftmargin=2mm,xrightmargin=1mm]
Nous venons de discuter du choix de [sujet]. Rédige la fiche de décision en
Markdown : contexte, options comparées, critères, décision retenue, alternatives
écartées avec leur raison, niveau de preuve (I, II ou III), références (marque
[à vérifier] si nécessaire), risques restants. N'invente aucune référence
ni aucun chiffre.
\end{Verbatim}

\CredixCapture{Un extrait du registre des décisions}{../images/guides/02-03_registre_decisions.png}

\emph{Figure 3. Un extrait du registre des décisions méthodologiques.}

\begin{center}\rule{0.5\linewidth}{0.5pt}\end{center}

\subsection{9. Étape 4 : choisir les langages et les
technologies}\label{9-uxe9tape-4--choisir-les-langages-et-les-technologies}

\textbf{À faire avant tout prototype.} Un mauvais choix technique se
paie à chaque étape suivante.

\textbf{Prompt 7 : proposer une pile technologique justifiée}

\begin{Verbatim}[breaklines=true,breakanywhere=true,fontsize=\footnotesize,frame=leftline,framerule=1.6pt,rulecolor=\color{credixblue},framesep=3mm,xleftmargin=2mm,xrightmargin=1mm]
À partir des exigences du cahier des charges (EF et ENF), propose une pile
technologique : langage(s), cadre de développement, base de données, outils
de test, moyen de lancement. Pour chaque choix : l'exigence qu'il sert,
deux alternatives écartées et pourquoi, les risques. Contraintes : je débute
sur [technologie], je dois pouvoir tout lancer en local en quelques commandes.
\end{Verbatim}

\textbf{Règle utilisée pour CREDIX :} \emph{chaque technologie répond à
une exigence formulée plus haut.} Exemples du mémoire : \textbf{Python}
pour l\textquotesingle apprentissage \textbf{et}
l\textquotesingle exploitation, afin d\textquotesingle éviter tout écart
silencieux entre le modèle entraîné et le modèle exécuté ;
\textbf{FastAPI} pour la vérification déclarative des requêtes et le
contrôle centralisé des droits ; \textbf{Next.js} avec
\textbf{TypeScript}, dont le typage fait apparaître les écarts de format
avant l\textquotesingle exécution ; \textbf{MongoDB} pour conserver des
demandes de structures variables sans migration ; un \textbf{service PDF
séparé} pour les comptes rendus.

\begin{quote}
\textbf{Dans CREDIX (séance du 30 mai).} Les décisions
d\textquotesingle architecture ont été prises point par point : quatre
conteneurs Docker (API, base de données, suivi des modèles, service
PDF), huit collections de données, des seuils de décision configurables
depuis l\textquotesingle interface plutôt qu\textquotesingle écrits en
dur, une promotion de modèle soumise à validation humaine.
\end{quote}

\begin{center}\rule{0.5\linewidth}{0.5pt}\end{center}

\subsection{10. Étape 5 : le document de
conception}\label{10-uxe9tape-5--le-document-de-conception}

\subsubsection{10.1 Ce que contient le
document}\label{101-ce-que-contient-le-document}

\begin{enumerate}
\def\labelenumi{\arabic{enumi}.}
\tightlist
\item
  \textbf{Contexte, objectifs, périmètre} ;
\item
  \textbf{Acteurs} et \textbf{cas d\textquotesingle utilisation}
  (diagramme et descriptions) ;
\item
  \textbf{Modèle du domaine} : classes, associations, multiplicités ;
\item
  \textbf{Scénarios principaux} : diagrammes de séquence ;
\item
  \textbf{Architecture} : couches, composants, déploiement ;
\item
  \textbf{Données} : tables ou collections ;
\item
  \textbf{Interfaces} : routes de l\textquotesingle API, avec exemples
  de requêtes et de réponses ;
\item
  \textbf{Critères d\textquotesingle acceptation} et \textbf{plan de
  tests} ;
\item
  \textbf{Glossaire}.
\end{enumerate}

Pour que Claude Code puisse \textbf{lire} les diagrammes, on les écrit
en \textbf{texte} avec \textbf{Mermaid}, qui s\textquotesingle affiche
automatiquement sur GitHub. Un exemple complet est fourni dans le
\textbf{guide 3} (le MVP).

\textbf{Prompt 8 : rédiger la conception}

\begin{Verbatim}[breaklines=true,breakanywhere=true,fontsize=\footnotesize,frame=leftline,framerule=1.6pt,rulecolor=\color{credixblue},framesep=3mm,xleftmargin=2mm,xrightmargin=1mm]
À partir du cahier des charges et des décisions du projet, rédige le document
de conception en Markdown avec les neuf sections suivantes : [copier la liste
du paragraphe 10.1]. Diagrammes en Mermaid. Règles UML : pas d'attribut dont
le type est une autre classe du diagramme (utiliser une association, avec
multiplicités aux deux extrémités) ; les types énumérés portent le stéréotype
<<enumeration>>. Chaque diagramme est précédé d'une phrase d'introduction et
suivi d'une légende. Ne rien inventer qui ne découle pas des documents du projet.
\end{Verbatim}

\textbf{Prompt 9 : vérifier la cohérence entre documents}

\begin{Verbatim}[breaklines=true,breakanywhere=true,fontsize=\footnotesize,frame=leftline,framerule=1.6pt,rulecolor=\color{credixblue},framesep=3mm,xleftmargin=2mm,xrightmargin=1mm]
Compare le cahier des charges et le document de conception. Liste : les
exigences non couvertes par la conception, les éléments de conception sans
exigence, les noms qui diffèrent d'un document à l'autre, les contradictions.
\end{Verbatim}

\begin{quote}
\textbf{Dans CREDIX.} La conception suit les règles du cours de génie
logiciel (UML : classes, séquences, cas d\textquotesingle utilisation,
composants, déploiement), puis a été dessinée en schémas
(\texttt{docs/\allowbreak{}diagrammes/\allowbreak{}}).
\end{quote}

\begin{center}\rule{0.5\linewidth}{0.5pt}\end{center}

\subsection{11. Étapes 6 et 7 : Claude Code, du plan au prototype
testé}\label{11-uxe9tapes-6-et-7--claude-code-du-plan-au-prototype-testuxe9}

\subsubsection{11.1 Étape 6 : le mode Plan et la
TODO}\label{111-uxe9tape-6--le-mode-plan-et-la-todo}

On ouvre le projet dans VS Code, on \textbf{dépose le document de
conception} dans le dossier, et on passe Claude Code en \textbf{mode
Plan} (guide 1, section 8).

\textbf{Prompt 10 : produire le plan de réalisation}

\begin{Verbatim}[breaklines=true,breakanywhere=true,fontsize=\footnotesize,frame=leftline,framerule=1.6pt,rulecolor=\color{credixblue},framesep=3mm,xleftmargin=2mm,xrightmargin=1mm]
/plan Lis docs/conception.md et docs/cahier_des_charges.md.
Propose un plan de réalisation en modules ordonnés (du plus indépendant au plus
dépendant). Pour chaque module : son objectif, les fichiers à créer, les tests
à écrire, le critère qui dit qu'il est terminé, et les dépendances. Ne crée
aucun fichier pour l'instant. Signale les points du document de conception
qui te paraissent ambigus.
\end{Verbatim}

Vous \textbf{relisez le plan}, l\textquotesingle annotez dans VS Code,
le faites corriger, puis vous l\textquotesingle{}\textbf{approuvez}. Ce
plan devient la \textbf{TODO} du projet.

\CredixCapture{Claude Code en mode Plan : plan de réalisation}{../images/guides/02-04_plan_todo.png}

\emph{Figure 4. Le plan produit à partir du document de conception.}

\subsubsection{11.2 Étape 7 : un module à la fois, avec des
tests}\label{112-uxe9tape-7--un-module-uxe0-la-fois-avec-des-tests}

\textbf{Prompt 11 : réaliser un module}

\begin{Verbatim}[breaklines=true,breakanywhere=true,fontsize=\footnotesize,frame=leftline,framerule=1.6pt,rulecolor=\color{credixblue},framesep=3mm,xleftmargin=2mm,xrightmargin=1mm]
Réalise le module [M2] du plan. Écris d'abord les tests, puis le code.
Lance les tests et montre-moi le résultat. Ne touche à aucun autre module.
Si tu dois modifier un fichier existant, dis-le avant.
\end{Verbatim}

\textbf{Prompt 12 : quand un test échoue}

\begin{Verbatim}[breaklines=true,breakanywhere=true,fontsize=\footnotesize,frame=leftline,framerule=1.6pt,rulecolor=\color{credixblue},framesep=3mm,xleftmargin=2mm,xrightmargin=1mm]
Le test [nom] échoue avec ce message : [coller le message]. Avant de modifier
quoi que ce soit, explique la cause probable et propose la correction la plus
petite possible.
\end{Verbatim}

Chaque module est \textbf{validé avant de passer au suivant}.

\subsubsection{11.3 Confier une mission à Claude Code : un exemple
réel}\label{113-confier-une-mission-uxe0-claude-code--un-exemple-ruxe9el}

Le prompt fixe \textbf{la mission, les règles et les livrables}. Voici
sa structure, dans le prompt réellement utilisé pour la phase 2 de
CREDIX.

\textbf{Prompt (extrait réel)}

\begin{Verbatim}[breaklines=true,breakanywhere=true,fontsize=\footnotesize,frame=leftline,framerule=1.6pt,rulecolor=\color{credixblue},framesep=3mm,xleftmargin=2mm,xrightmargin=1mm]
Tu travailles sur mon mémoire CREDIX (scoring de crédit). Le notebook est dans
le workspace. Lis d'abord EN ENTIER le fichier HANDOVER_CLAUDE_CODE_PHASE2.md
(contexte, règles absolues, variables disponibles) avant de coder.

Ta mission : écrire les cellules de la Section 12.3+ qui démontrent [...].
Ce n'est pas un concours d'outils.

RÈGLES
- On étend, on ne supprime pas. Nouvelles cellules en Section 12.3+ UNIQUEMENT.
- Tu n'exécutes pas le notebook. Tu testes ta logique en local sur données
  factices, puis tu livres le code prêt à coller.
- Anti-fuite : tout ce qui s'apprend s'apprend sur le TRAIN.
\end{Verbatim}

Les \textbf{règles absolues} du document de passation associé : \emph{«
On étend, on ne supprime pas »} ; \emph{« Aucun chiffre inventé : le
code produit les chiffres ; on ne les anticipe jamais dans les
commentaires ou les rapports »} ; et la \textbf{division du travail} :
\textbf{Claude Code écrit et teste sur des données factices ;
l\textquotesingle auteur exécute sur les vraies données et renvoie les
sorties.}

\begin{center}\rule{0.5\linewidth}{0.5pt}\end{center}

\subsection{12. Étape 8 : livrer, lancer, tester à la
main}\label{12-uxe9tape-8--livrer-lancer-tester-uxe0-la-main}

\textbf{Prompt 13 : préparer la livraison et le README}

\begin{Verbatim}[breaklines=true,breakanywhere=true,fontsize=\footnotesize,frame=leftline,framerule=1.6pt,rulecolor=\color{credixblue},framesep=3mm,xleftmargin=2mm,xrightmargin=1mm]
Prépare la livraison du projet. 1) Vérifie que tous les tests passent.
2) Rédige le README de lancement pour un débutant : prérequis, commandes
exactes dans l'ordre, ce qu'on doit voir après chaque commande, et que faire
en cas d'erreur. 3) Génère une archive ZIP du projet, sans les fichiers
secrets ni les dossiers volumineux.
\end{Verbatim}

Puis, \textbf{vous} : dézippez dans un \textbf{dossier vide}, suivez le
README \textbf{à la lettre}, notez \textbf{tout ce qui ne marche pas},
et renvoyez ces constats à Claude.

\begin{quote}
\textbf{Dans CREDIX (séance du 30 mai).} Le backend a été livré sous
forme d\textquotesingle une archive de 51 fichiers, avec
l\textquotesingle{}\textbf{ordre de lancement exact} : copier le fichier
de configuration, copier les fichiers du modèle, lancer les conteneurs,
créer l\textquotesingle administrateur, générer les phrases
d\textquotesingle explication, enregistrer le modèle, charger les
clients de démonstration, vérifier le point de santé.

\textbf{Pour ce dossier de remise}, le test du README sur une copie
neuve a révélé que la \textbf{compilation de production du frontend
n\textquotesingle avait jamais abouti} : quatre erreurs de types que le
mode développement ne signalait pas. \textbf{Ce qui n\textquotesingle a
jamais été testé de bout en bout n\textquotesingle est pas testé.}
\end{quote}

\begin{center}\rule{0.5\linewidth}{0.5pt}\end{center}

\subsection{13. Étape 9 : modifier, documenter,
transmettre}\label{13-uxe9tape-9--modifier-documenter-transmettre}

\textbf{Prompt 14 : ajouter une fonction (toujours en mode Plan)}

\begin{Verbatim}[breaklines=true,breakanywhere=true,fontsize=\footnotesize,frame=leftline,framerule=1.6pt,rulecolor=\color{credixblue},framesep=3mm,xleftmargin=2mm,xrightmargin=1mm]
/plan Je veux ajouter [fonction]. Analyse l'impact : quels fichiers sont
concernés, quels tests existants risquent de casser, quelles nouvelles
données sont nécessaires. Propose les étapes et les tests. Ne modifie rien
avant mon accord.
\end{Verbatim}

Puis on boucle sur le cycle de session (section 3) : le travail est
consigné, la passation écrite.

\begin{quote}
\textbf{Dans CREDIX.} Le frontend d\textquotesingle origine avait été
généré à partir de maquettes. Il a été reconstruit (\texttt{credix-v2})
avec Claude Code : cahier des charges de l\textquotesingle interface,
système de design, puis réalisation \textbf{par phases} (fondations,
espace Agent, espace Superviseur, espace Administrateur), \textbf{chaque
phase vérifiée avant la suivante}.
\end{quote}

\begin{center}\rule{0.5\linewidth}{0.5pt}\end{center}

\subsection{14. Les erreurs et les pièges
réels}\label{14-les-erreurs-et-les-piuxe8ges-ruxe9els}

\textbf{Claude se trompe, comme un humain.} La méthode ne
l\textquotesingle empêche pas : elle \textbf{permet de le voir}. Voici
des erreurs réelles de CREDIX, et la règle qui en est sortie.

\begin{longtable}[]{@{}
  >{\raggedright\arraybackslash}p{(\linewidth - 4\tabcolsep) * \real{0.3233}}
  >{\raggedright\arraybackslash}p{(\linewidth - 4\tabcolsep) * \real{0.3233}}
  >{\raggedright\arraybackslash}p{(\linewidth - 4\tabcolsep) * \real{0.3233}}@{}}
\toprule\noalign{}
\begin{minipage}[b]{\linewidth}\raggedright
Ce qui s\textquotesingle est passé
\end{minipage} & \begin{minipage}[b]{\linewidth}\raggedright
Comment c\textquotesingle est apparu
\end{minipage} & \begin{minipage}[b]{\linewidth}\raggedright
Règle qui en découle
\end{minipage} \\
\midrule\noalign{}
\endhead
\bottomrule\noalign{}
\endlastfoot
Un seuil de détection d\textquotesingle anomalie allait être calculé sur
« les bons payeurs », c\textquotesingle est-à-dire avec la variable
cible, dans un mécanisme \textbf{non supervisé} (fuite
d\textquotesingle information) & L\textquotesingle auteur a relevé
l\textquotesingle erreur ; les seuils ont été \textbf{chargés depuis les
fichiers du modèle}, non recalculés & \textbf{Relire} chaque proposition
avec ses contraintes de rigueur en tête \\
Un argument disait que LightGBM « gagnait nettement sur le rappel » &
Une comparaison mesurée a montré des écarts non significatifs entre
LightGBM et XGBoost ; \textbf{l\textquotesingle argument a été rejeté},
LightGBM gardé sur d\textquotesingle autres critères & Une décision «
définitive » peut être révisée si un chiffre la contredit : \textbf{le
registre garde la trace} \\
Des documents externes donnaient un indicateur de 0,13 pour une variable
& Vérifié contre l\textquotesingle exécution réelle : \textbf{0,0199} &
\textbf{Vérifier tout chiffre reçu d\textquotesingle une source
externe}, même si le raisonnement est juste \\
Un défaut d\textquotesingle alignement dans le calcul de
l\textquotesingle IV (section 8) & Contrôle cellule par cellule &
\textbf{Vérifier les résultats étonnants}, y compris les résultats «
trop faibles » \\
Un test comparait « défaut des accordés avec anomalie » et « accordés
sains » : hors sujet, car une anomalie n\textquotesingle est pas un
défaut & Recentrage acté avec l\textquotesingle auteur & Demander :
\emph{« ce test répond-il vraiment à mon objectif ? »} \\
Un environnement installait une version trop ancienne
d\textquotesingle une bibliothèque & Message d\textquotesingle erreur
explicite & \textbf{Figer les versions} des bibliothèques \\
La compilation de production du frontend n\textquotesingle avait jamais
été lancée & Test sur copie neuve (ce dossier) & \textbf{Tester comme le
fera l\textquotesingle utilisateur final} \\
\end{longtable}

\textbf{Les réflexes qui en découlent :} relire ; mesurer ; garder la
trace ; refuser les chiffres non sourcés ; ne jamais présenter le
système autrement qu\textquotesingle il n\textquotesingle est ; tester
sur une machine neuve.

\begin{center}\rule{0.5\linewidth}{0.5pt}\end{center}

\subsection{15. Bibliothèque de
prompts}\label{15-bibliothuxe8que-de-prompts}

\begin{longtable}[]{@{}
  >{\raggedright\arraybackslash}p{(\linewidth - 4\tabcolsep) * \real{0.0909}}
  >{\raggedright\arraybackslash}p{(\linewidth - 4\tabcolsep) * \real{0.1718}}
  >{\raggedright\arraybackslash}p{(\linewidth - 4\tabcolsep) * \real{0.7073}}@{}}
\toprule\noalign{}
\begin{minipage}[b]{\linewidth}\raggedright
N°
\end{minipage} & \begin{minipage}[b]{\linewidth}\raggedright
Moment
\end{minipage} & \begin{minipage}[b]{\linewidth}\raggedright
Objet
\end{minipage} \\
\midrule\noalign{}
\endhead
\bottomrule\noalign{}
\endlastfoot
\textbf{S} & Début de session & Prompt de démarrage : ordre de lecture,
attendre le feu vert (section 3.1) \\
\textbf{F} & Fin de session & Rapport de passation (section 3.3) \\
\textbf{0} & Étape 0 & Faire l\textquotesingle inventaire du projet \\
\textbf{1} & Étape 1 & Cadrer le besoin en se faisant interroger \\
\textbf{2} & Étape 1 & Rédiger le cahier des charges \\
\textbf{3} & Étape 1 & Faire critiquer le cahier des charges \\
\textbf{4} & Étape 2 & Balayer la littérature (recherche approfondie) \\
\textbf{5} & Étape 2 & Interroger un article chargé \\
\textbf{6} & Étape 3 & Produire une fiche de décision \\
\textbf{7} & Étape 4 & Proposer une pile technologique \\
\textbf{8} & Étape 5 & Rédiger le document de conception \\
\textbf{9} & Étape 5 & Vérifier la cohérence entre documents \\
\textbf{10} & Étape 6 & Produire le plan de réalisation (mode Plan) \\
\textbf{11} & Étape 7 & Réaliser un module avec tests \\
\textbf{12} & Étape 7 & Traiter un test qui échoue \\
\textbf{13} & Étape 8 & Préparer la livraison et le README \\
\textbf{14} & Étape 9 & Ajouter ou modifier une fonction (mode Plan) \\
\end{longtable}

Deux prompts utiles à toutes les étapes :

\textbf{Prompt 15 : l\textquotesingle avocat du diable}

\begin{Verbatim}[breaklines=true,breakanywhere=true,fontsize=\footnotesize,frame=leftline,framerule=1.6pt,rulecolor=\color{credixblue},framesep=3mm,xleftmargin=2mm,xrightmargin=1mm]
Donne-moi les cinq meilleures objections à [ce choix, ce plan, ce texte].
Pour chacune : est-elle décisive ? comment la traiter ?
\end{Verbatim}

\textbf{Prompt 16 : la chasse aux chiffres non sourcés}

\begin{Verbatim}[breaklines=true,breakanywhere=true,fontsize=\footnotesize,frame=leftline,framerule=1.6pt,rulecolor=\color{credixblue},framesep=3mm,xleftmargin=2mm,xrightmargin=1mm]
Relis ce document. Liste chaque chiffre, chaque affirmation factuelle et chaque
référence. Pour chacun : d'où vient-il (document du projet, calcul, littérature,
ou nulle part) ? Signale tout ce qui n'a pas de source.
\end{Verbatim}

\begin{center}\rule{0.5\linewidth}{0.5pt}\end{center}

\subsection{16. Modèles de documents}\label{16-moduxe8les-de-documents}

\subsubsection{16.1 Fiche de décision}\label{161-fiche-de-duxe9cision}

\begin{Verbatim}[breaklines=true,breakanywhere=true,fontsize=\footnotesize,frame=leftline,framerule=1.6pt,rulecolor=\color{credixblue},framesep=3mm,xleftmargin=2mm,xrightmargin=1mm]
## Décision D-[n] : [titre]
- **Date :** [jj/mm/aaaa]
- **Contexte :** [le problème, les contraintes]
- **Options comparées :** [A], [B], [C]
- **Critères :** [explicabilité, coût, stabilité, ...]
- **Décision :** [option retenue]
- **Pourquoi :** [justification en quelques lignes]
- **Alternatives écartées :** [option : raison]
- **Niveau de preuve :** [I démontré / II littérature / III hypothèse assumée]
- **Références :** [auteur, année] [à vérifier si besoin]
- **Risques restants :** [...]
- **Validée par :** [nom] le [date]
\end{Verbatim}

\subsubsection{16.2 Entrée du journal de bord (ajout
seulement)}\label{162-entruxe9e-du-journal-de-bord-ajout-seulement}

\begin{Verbatim}[breaklines=true,breakanywhere=true,fontsize=\footnotesize,frame=leftline,framerule=1.6pt,rulecolor=\color{credixblue},framesep=3mm,xleftmargin=2mm,xrightmargin=1mm]
## Session [n] : [date]
**Objectif :** [...]
**Réalisé :** [liste]
**Décisions prises :** [liste, avec renvoi vers le registre]
**En attente / à faire à la prochaine session :** [liste ordonnée]
\end{Verbatim}

\subsubsection{16.3 Instructions du projet (six
sections)}\label{163-instructions-du-projet-six-sections}

\begin{Verbatim}[breaklines=true,breakanywhere=true,fontsize=\footnotesize,frame=leftline,framerule=1.6pt,rulecolor=\color{credixblue},framesep=3mm,xleftmargin=2mm,xrightmargin=1mm]
# Instructions du projet Claude : [titre du projet]
## Section 1 : Contexte
## Section 2 : Ressources permanentes (liste classée, avec « consulter pour... »)
## Section 3 : Décisions techniques figées
## Section 4 : Consignes permanentes (langue, niveau, style, à faire, à ne pas faire)
## Section 5 : Workflow de collaboration (début / pendant / fin de session)
## Section 6 : État courant (mis à jour à chaque session)
\end{Verbatim}

\subsubsection{16.4 Document de
passation}\label{164-document-de-passation}

\begin{Verbatim}[breaklines=true,breakanywhere=true,fontsize=\footnotesize,frame=leftline,framerule=1.6pt,rulecolor=\color{credixblue},framesep=3mm,xleftmargin=2mm,xrightmargin=1mm]
# Passation : [projet], [phase]
## 1. Contexte minimal
## 2. Règles absolues (non négociables)
## 3. Où en est le projet (ce qui est fait : ne pas refaire)
## 4. Ce qui reste (le travail demandé)
## 5. Fichiers, variables, commandes disponibles
\end{Verbatim}

\subsubsection{16.5 Liste de pièges}\label{165-liste-de-piuxe8ges}

\begin{Verbatim}[breaklines=true,breakanywhere=true,fontsize=\footnotesize,frame=leftline,framerule=1.6pt,rulecolor=\color{credixblue},framesep=3mm,xleftmargin=2mm,xrightmargin=1mm]
## Pièges déjà identifiés : à ne jamais reproduire
- **[Sujet].** [Ce qui s'est passé] → [ce qu'il faut faire à la place].
\end{Verbatim}

\begin{center}\rule{0.5\linewidth}{0.5pt}\end{center}

\subsection{17. Rejouer la méthode : le
mini-exercice}\label{17-rejouer-la-muxe9thode--le-mini-exercice}

Pour vérifier que la méthode fonctionne \textbf{sans connaître CREDIX},
le \textbf{guide 3} fournit un document de conception réduit (MVP)
d\textquotesingle un mini-système de scoring de crédit. Le protocole de
test devant l\textquotesingle encadrant est dans la \textbf{fiche de
démonstration}. En résumé :

\begin{enumerate}
\def\labelenumi{\arabic{enumi}.}
\tightlist
\item
  Ouvrir un dossier vide dans VS Code et y déposer
  \texttt{03\_\allowbreak{}conception\_\allowbreak{}MVP.\allowbreak{}md}.
\item
  Ouvrir Claude Code, passer en \textbf{mode Plan}, et donner le
  \textbf{prompt 10}.
\item
  \textbf{Relire le plan} : est-il fidèle au document ? Les modules
  sont-ils dans le bon ordre ?
\item
  Approuver, réaliser \textbf{un module} avec le \textbf{prompt 11}, et
  vérifier que ses tests passent.
\end{enumerate}

Si le plan est cohérent et que le premier module passe ses tests,
\textbf{la méthode est validée}.

\end{document}
